%% kapa: exam=1 round=28 year=2017 period=2 form=A subject=appraisal_law subject_ko=감정평가관계법규 seq=11 number=11
%% source: Q-Net 감정평가사 기출문제 (cst00309) artlSeq=5199603 file=extracted/1008/5199603/2017년 제28회 감정평가사 1차시험 문제지/2교시 A형.pdf
%% conversion: kordoc 4.24.51 -> tools/gen1.py (auto); review=proofread-sonnet; flags=none
\begin{question}{11}
국토의 계획 및 이용에 관한 법령상 용도지역에서의 용적률의 최대한도가 큰 순서대로 나열된 것은? (단, 조례는 고려하지 않음)
\begin{framedtext}
\noindent\begin{tabular}{@{}p{0.42\linewidth}p{0.42\linewidth}@{}}
ㄱ. 제1종일반주거지역 & ㄴ. 일반상업지역 \\
ㄷ. 전용공업지역 & ㄹ. 자연녹지지역 \\
\end{tabular}
\end{framedtext}
\begin{choices}
  \item ㄱ - ㄴ - ㄷ - ㄹ
  \item ㄴ - ㄱ - ㄷ - ㄹ
  \item ㄴ - ㄷ - ㄱ - ㄹ
  \item ㄷ - ㄱ - ㄴ - ㄹ
  \item ㄷ - ㄴ - ㄱ - ㄹ
\end{choices}
\end{question}
\answer{3}
